\documentclass[10pt,a4paper]{amsart}
\usepackage{amsmath,amssymb,amsthm,graphicx}
% Retained preamble for the Stacks Project chapter divisors.tex (Divisors).
% The chapter file pinned-original-source.tex carries only \input{preamble} on its first line, so no
% preamble text lives inside the chapter itself.  The Stacks Project uses one single shared preamble for
% every chapter; the snapshot stages it verbatim in sources/stacks-chow/preamble.tex.  The block below is
% copied mechanically, line for line, from that file's lines 14-18 (the Xy-pic setup) and lines 154-217
% (the native theorem declarations together with the mathematical macros).  Nothing was added, removed or
% redefined.  The wrapper class and the amsmath/amssymb/amsthm/graphicx base packages are supplied by the
% assembler, so the class/\externaldocument part of the shared preamble is not repeated here.
% For commutative diagrams we use Xy-pic
\usepackage[all]{xy}

% We use 2cell for 2-commutative diagrams.
\xyoption{2cell}
% Theorem environments.
%
\theoremstyle{plain}
\newtheorem{theorem}[subsection]{Theorem}
\newtheorem{proposition}[subsection]{Proposition}
\newtheorem{lemma}[subsection]{Lemma}

\theoremstyle{definition}
\newtheorem{definition}[subsection]{Definition}
\newtheorem{example}[subsection]{Example}
\newtheorem{exercise}[subsection]{Exercise}
\newtheorem{situation}[subsection]{Situation}

\theoremstyle{remark}
\newtheorem{remark}[subsection]{Remark}
\newtheorem{remarks}[subsection]{Remarks}

\numberwithin{equation}{subsection}

% Macros
%
\def\lim{\mathop{\mathrm{lim}}\nolimits}
\def\colim{\mathop{\mathrm{colim}}\nolimits}
\def\Spec{\mathop{\mathrm{Spec}}}
\def\Hom{\mathop{\mathrm{Hom}}\nolimits}
\def\Ext{\mathop{\mathrm{Ext}}\nolimits}
\def\SheafHom{\mathop{\mathcal{H}\!\mathit{om}}\nolimits}
\def\SheafExt{\mathop{\mathcal{E}\!\mathit{xt}}\nolimits}
\def\Sch{\mathit{Sch}}
\def\Mor{\mathop{\mathrm{Mor}}\nolimits}
\def\Ob{\mathop{\mathrm{Ob}}\nolimits}
\def\Sh{\mathop{\mathit{Sh}}\nolimits}
\def\NL{\mathop{N\!L}\nolimits}
\def\CH{\mathop{\mathrm{CH}}\nolimits}
\def\proetale{{pro\text{-}\acute{e}tale}}
\def\etale{{\acute{e}tale}}
\def\QCoh{\mathit{QCoh}}
\def\Ker{\mathop{\mathrm{Ker}}}
\def\Im{\mathop{\mathrm{Im}}}
\def\Coker{\mathop{\mathrm{Coker}}}
\def\Coim{\mathop{\mathrm{Coim}}}

% Boxtimes
%
\DeclareMathSymbol{\boxtimes}{\mathbin}{AMSa}{"02}

%
% Macros for moduli stacks/spaces
%
\def\QCohstack{\mathcal{QC}\!\mathit{oh}}
\def\Cohstack{\mathcal{C}\!\mathit{oh}}
\def\Spacesstack{\mathcal{S}\!\mathit{paces}}
\def\Quotfunctor{\mathrm{Quot}}
\def\Hilbfunctor{\mathrm{Hilb}}
\def\Curvesstack{\mathcal{C}\!\mathit{urves}}
\def\Polarizedstack{\mathcal{P}\!\mathit{olarized}}
\def\Complexesstack{\mathcal{C}\!\mathit{omplexes}}
% \Pic is the operator that assigns to X its picard group, usage \Pic(X)
% \Picardstack_{X/B} denotes the Picard stack of X over B
% \Picardfunctor_{X/B} denotes the Picard functor of X over B
\def\Pic{\mathop{\mathrm{Pic}}\nolimits}
\def\Picardstack{\mathcal{P}\!\mathit{ic}}
\def\Picardfunctor{\mathrm{Pic}}
\def\Deformationcategory{\mathcal{D}\!\mathit{ef}}

% --- appended 2026-10-12 (Chief): the project own definitions of the reference/slogan/history
% display environments, copied mechanically from the shared Stacks preamble lines 8-12
% (\usepackage{verbatim} plus the three \newenvironment lines, which the chapter preamble keeps
% verbatim from the project preamble). Nothing added, removed or redefined: these environments
% carry no mathematical content and the project itself typesets them as comments. Without these
% three lines a counted statement whose own text uses them fails with Environment slogan undefined.
\usepackage{verbatim}
\newenvironment{reference}{\comment}{\endcomment}
\newenvironment{slogan}{\comment}{\endcomment}
\newenvironment{history}{\comment}{\endcomment}
\usepackage{cleveref}
\crefname{originallemma}{original lemma}{original lemmas}
\Crefname{originallemma}{Original lemma}{Original lemmas}
\makeatletter
\newlabel{lemma-locally-free-rank-r-pullback}{{05P8}{source}{}{sourcepointer-1}{}}
\newlabel{lemma-locally-free-rank-r-pullback@cref}{{[originallemma][0][]05P8}{[1][0][]source}{}{}{}}
\makeatother
\crefname{originallemma}{original lemma}{original lemmas}
\Crefname{originallemma}{Original lemma}{Original lemmas}
\makeatletter
\newlabel{properties-lemma-finite-locally-free}{{05P2}{source}{}{sourcepointer-2}{}}
\newlabel{properties-lemma-finite-locally-free@cref}{{[originallemma][0][]05P2}{[1][0][]source}{}{}{}}
\makeatother
\title{The locus where a finitely presented module is flat is a disjoint union of rank strata}
\author{The Stacks Project Authors}
\date{stacks-project:divisors.tex:local-snapshot-20261012; GNU FDL 1.2 or later}
\begin{document}
\maketitle
\setcounter{subsection}{1}
\noindent\textit{Editorial note.} All delivered mathematical text is the author's own, copied from the pinned Stacks Project chapter file divisors.tex; no statement or proof step was rewritten, reordered or removed.\par\medskip
\noindent\textit{Editorial note.} Carrier window: statement and proof both lie inside the window 0-2000 granted to this lane by CONTROL 8.83.\par\medskip
\noindent\textit{Editorial note.} The card is taken because the representability of F\_flat is proved in full; the trailing sentence 'The other statements also follow from the lemma' delegates only the 'moreover' to the cited rank lemma, and that delegation is recorded here for the reader.\par\medskip
\noindent\textit{Editorial note.} The two references of the counted spans are recorded as non-load-bearing navigation pointers; the argument that consumes them stays inside the counted proof span.\par\medskip
\noindent\textit{Editorial note.} Preamble: sources/stacks-divisors/retained-preamble.tex (88 lines, verbatim present).\par\medskip
\noindent\textit{Allowed standard prerequisites (editorial).}\begin{itemize}
\item The closed subschemes and functors measuring the locally free rank. Lemma lemma-locally-free-rank-r-pullback, used at proof lines 1466-1468 and 1491-1499. Supplies Z\_r and the functors F\_r representing the rank-r locus.
\item A flat module of finite presentation is finite locally free. Properties, Lemma properties-lemma-finite-locally-free, used at proof lines 1487-1489. Turns flatness into a coproduct decomposition by rank.
\end{itemize}\medskip
\ifdefined\hypertarget\hypertarget{sourcepointer-1}{}\fi
\noindent\textit{Source reference.} Original lemma 05P8: focused-seeds/AGP/sources/stacks-snapshot-20261011/tags.txt, source lines 7332--7332. Navigation pointer into the Stacks Project; the standard fact it supplies is applied inside the counted proof span and the argument resting on it stays there.\par\medskip
\ifdefined\hypertarget\hypertarget{sourcepointer-2}{}\fi
\noindent\textit{Source reference.} Original lemma 05P2: focused-seeds/AGP/sources/stacks-snapshot-20261011/tags.txt, source lines 7326--7326. Navigation pointer into the Stacks Project; the standard fact it supplies is applied inside the counted proof span and the argument resting on it stays there.\par\medskip
\begin{lemma}
\label{lemma-finite-presentation-module}
Let $S$ be a scheme. Let $\mathcal{F}$ be an $\mathcal{O}_S$-module
of finite presentation. Let $S = Z_{-1} \supset Z_0 \supset Z_1 \supset \ldots$
be as in Lemma \ref{lemma-locally-free-rank-r-pullback}.
Set $S_r = Z_{r - 1} \setminus Z_r$.
Then $S' = \coprod_{r \geq 0} S_r$ represents the functor
$$
F_{flat} : \Sch/S \longrightarrow \textit{Sets},\quad\quad
T \longmapsto
\left\{
\begin{matrix}
\{*\} & \text{if }\mathcal{F}_T\text{ flat over }T\\
\emptyset & \text{otherwise}
\end{matrix}
\right.
$$
Moreover, $\mathcal{F}|_{S_r}$ is locally free of rank $r$ and the
morphisms $S_r \to S$ and $S' \to S$ are of finite presentation.
\end{lemma}

\begin{proof}
Suppose that $g : T \to S$ is a morphism of schemes such that the pullback
$\mathcal{F}_T = g^*\mathcal{F}$ is flat. Then $\mathcal{F}_T$ is a flat
$\mathcal{O}_T$-module of finite presentation. Hence
$\mathcal{F}_T$ is finite locally free, see
Properties, Lemma \ref{properties-lemma-finite-locally-free}.
Thus $T = \coprod_{r \geq 0} T_r$, where $\mathcal{F}_T|_{T_r}$ is locally
free of rank $r$. This implies that
$$
F_{flat} = \coprod\nolimits_{r \geq 0} F_r
$$
in the category of Zariski sheaves on $\Sch/S$ where $F_r$ is as in
Lemma \ref{lemma-locally-free-rank-r-pullback}. It follows
that $F_{flat}$ is represented by
$\coprod_{r \geq 0} (Z_{r - 1} \setminus Z_r)$ where
$Z_r$ is as in
Lemma \ref{lemma-locally-free-rank-r-pullback}.
The other statements also follow from the lemma.
\end{proof}

\end{document}
